% GPU CrocSort: Shared-Memory K-Way Merge Trees for GPU External Sort
% LaTeX source - IEEE Conference Format
%
% Compile with: pdflatex gpu_crocsort.tex && bibtex gpu_crocsort && pdflatex gpu_crocsort.tex && pdflatex gpu_crocsort.tex

\documentclass[conference,10pt]{IEEEtran}

% ─── Packages ────────────────────────────────────────────────────────────────
\usepackage[utf8]{inputenc}
\usepackage[T1]{fontenc}
\usepackage{amsmath,amssymb,amsfonts}
\usepackage{graphicx}
\usepackage{booktabs}
\usepackage{multirow}
\usepackage{xcolor}
\usepackage{hyperref}
\usepackage{cleveref}
\usepackage{microtype}
\usepackage{subcaption}
\usepackage{enumitem}
\usepackage{balance}
\usepackage{xspace}

% Algorithm packages
\usepackage[ruled,vlined,linesnumbered]{algorithm2e}
\SetAlgoLined
\SetKwInput{KwInput}{Input}
\SetKwInput{KwOutput}{Output}
\SetKwProg{Fn}{Function}{:}{}
\SetKwProg{Kernel}{Kernel}{:}{}
\SetKwComment{Comment}{// }{}

% ─── Custom Commands ─────────────────────────────────────────────────────────
\newcommand{\crocsort}{\textsc{CrocSort}\xspace}
\newcommand{\gpucrocsort}{\textsc{GPU CrocSort}\xspace}
\newcommand{\ovc}{\textsc{OVC}\xspace}
\newcommand{\hmem}{\textsf{HBM}\xspace}
\newcommand{\smem}{\textsf{SMEM}\xspace}
\newcommand{\BigO}[1]{\mathcal{O}\!\left(#1\right)}
\newcommand{\todo}[1]{\textcolor{red}{\textbf{TODO:} #1}}

% Placeholder figure command
\newcommand{\placeholderfig}[2]{%
  \begin{center}
    \fbox{\parbox{0.9\columnwidth}{\centering\vspace{#1}\textit{[Figure placeholder: #2]}\vspace{#1}}}
  \end{center}
}

% ─── Document ─────────────────────────────────────────────────────────────────
\begin{document}

\title{GPU CrocSort: Shared-Memory K-Way Merge Trees\\for GPU External Sort}

\author{
  \IEEEauthorblockN{Draft}
  \IEEEauthorblockA{Working Paper\\
  \today}
}

\maketitle

% ══════════════════════════════════════════════════════════════════════════════
% ABSTRACT
% ══════════════════════════════════════════════════════════════════════════════
\begin{abstract}
External merge sort remains the method of choice for sorting datasets that
exceed available memory, yet GPU implementations face a fundamental tension
between merge fan-in and thread utilization.  A high fan-in (K-way merge)
reduces the number of merge passes and thus total memory traffic, but
traditional K-way merging serializes into a single active thread, squandering
the massive parallelism that GPUs provide.  Conversely, 2-way merge-path
decompositions keep every thread busy but demand $\lceil\log_2 N\rceil$
passes, multiplying bandwidth consumption.

We present \gpucrocsort, a CUDA-based external merge sort that resolves this
trade-off.  The key contribution is a \emph{shared-memory K-way merge tree}
that decomposes a $K$-way merge ($K{=}8$) into $\log_2 K{=}3$ levels of
fully parallel 2-way merge-path operations entirely within GPU shared memory.
Combined with Offset-Value Coding (\ovc) for single-instruction key
comparison, \gpucrocsort achieves the low pass count of K-way merging
\emph{and} the full thread utilization of 2-way merge-path---simultaneously.
On standard GenSort benchmarks, the 8-way strategy reduces HBM traffic by
$3\times$ compared to the 2-way baseline while maintaining full occupancy
across all streaming multiprocessors.
\end{abstract}

\begin{IEEEkeywords}
GPU sorting, external merge sort, CUDA, shared memory, K-way merge,
merge path, offset-value coding
\end{IEEEkeywords}

% ══════════════════════════════════════════════════════════════════════════════
% 1  INTRODUCTION
% ══════════════════════════════════════════════════════════════════════════════
\section{Introduction}
\label{sec:intro}

Sorting is among the most fundamental operations in data processing, and
external sorting---sorting datasets that exceed the capacity of fast
memory---remains indispensable in database systems, scientific computing, and
large-scale analytics.  As dataset sizes grow, the performance of external
sort is dominated by the number of \emph{merge passes}, each of which reads
and writes the entire dataset through the memory hierarchy.

Modern GPUs offer an order-of-magnitude advantage in memory bandwidth over
CPUs: NVIDIA's A100 delivers up to 2,039\,GB/s of HBM2e bandwidth compared
to roughly 200\,GB/s of DDR5 bandwidth on contemporary CPU platforms.  This
bandwidth advantage makes GPUs a natural target for merge-sort
acceleration---provided the merge algorithm can keep the thousands of GPU
threads productively occupied.

Herein lies the central dilemma.  The \emph{merge-path}
decomposition~\cite{green2012} enables a 2-way merge to be executed in
parallel across all GPU threads, but 2-way merging requires
$\lceil\log_2 N\rceil$ passes over the data, where $N$ is the number of
initial sorted runs.  For $N{=}19{,}531$ runs (a 10\,GB dataset with
512-record blocks), this amounts to 15 passes and 30\,GB of total HBM
traffic.  K-way merging with $K{=}8$ requires only
$\lceil\log_8 N\rceil = 5$ passes and 10\,GB of traffic---a $3\times$
reduction---but traditional K-way merge implementations employ a priority
queue (loser tree) that serializes into a single active thread per merge,
leaving the vast majority of GPU cores idle.

\textbf{Contributions.}\quad We present \gpucrocsort, which resolves this
tension through three contributions:

\begin{enumerate}[leftmargin=*,topsep=2pt,itemsep=1pt]
  \item \textbf{Shared-memory K-way merge tree.}  We decompose a $K$-way
    merge into $\log_2 K$ levels of 2-way merge-path operations that
    execute entirely in shared memory.  Each level is fully parallel (all
    threads active), and only the final level writes to HBM, achieving both
    high fan-in \emph{and} full thread utilization.
  \item \textbf{GPU-adapted OVC encoding.}  We port the Offset-Value Coding
    scheme from the CPU \crocsort~\cite{crocsort2026} to CUDA, exploiting
    the 32-bit packed format for single-instruction comparison in GPU
    registers.
  \item \textbf{End-to-end GPU external sort.}  We integrate run generation
    via bitonic sort, multi-pass merge, and a streaming external-sort
    pipeline with pinned host memory and asynchronous transfers.
\end{enumerate}

The remainder of this paper is organized as follows.
\Cref{sec:background} surveys background and related work.
\Cref{sec:architecture} presents the overall system architecture.
\Cref{sec:kway} details the shared-memory K-way merge tree.
\Cref{sec:ovc} describes GPU-adapted OVC encoding.
\Cref{sec:analysis} provides theoretical analysis.
\Cref{sec:experiments} covers experimental setup and results.
\Cref{sec:related} discusses related work in depth.
\Cref{sec:future} outlines future directions, and
\Cref{sec:conclusion} concludes.


% ══════════════════════════════════════════════════════════════════════════════
% 2  BACKGROUND
% ══════════════════════════════════════════════════════════════════════════════
\section{Background}
\label{sec:background}

\subsection{External Merge Sort}
External merge sort proceeds in two phases.  In the \emph{run generation}
phase, chunks of data that fit in fast memory are sorted to produce initial
sorted runs.  In the \emph{merge} phase, runs are repeatedly merged until a
single sorted output remains.  The number of passes through the data is
determined by the \emph{merge fan-in}~$K$: for $N$ initial runs, the merge
phase requires $\lceil\log_K N\rceil$ passes.

\subsection{CrocSort and Offset-Value Coding}
\crocsort~\cite{crocsort2026}, presented at VLDB 2026, introduced an
efficient CPU external merge sort built on \emph{Offset-Value Coding}
(\ovc)~\cite{ovc_original}.  An OVC is a 32-bit packed integer that encodes
the comparison result between a record's key and a reference key:

\begin{equation}
  \text{OVC}_{32} = \underbrace{\text{flag}}_{3\text{ bits}} \;\|\;
                     \underbrace{\text{arity\_minus\_offset}}_{13\text{ bits}} \;\|\;
                     \underbrace{\text{value}}_{16\text{ bits}}
  \label{eq:ovc_format}
\end{equation}

The flag field distinguishes five record types (\texttt{EarlyFence},
\texttt{Duplicate}, \texttt{Normal}, \texttt{Initial}, \texttt{LateFence}),
chosen so that the natural \texttt{uint32\_t} ordering of the packed OVC
matches the desired sort order.  This enables single-instruction comparison:
comparing two OVC values as unsigned 32-bit integers yields the correct
merge ordering without accessing the underlying variable-length key.

\crocsort also introduced an analytical model based on the \emph{M/D ratio}
(memory size $M$ divided by page size $D$) that governs the optimal merge
fan-in and predicts I/O costs.  We adapt this model to the GPU setting,
substituting HBM capacity for $M$ and streaming-block size for $D$.

\subsection{GPU Merge Primitives}
\label{sec:bg_gpu_merge}

\paragraph{Merge Path.}
Green et al.~\cite{green2012} introduced the \emph{merge-path}
decomposition, which frames a 2-way merge as a path through a
two-dimensional grid.  By partitioning the path among threads, every thread
computes an independent, equally-sized portion of the merged output.  This
guarantees load-balanced, fully parallel 2-way merging.

\paragraph{GPU Multiway Mergesort.}
Chatterjee et al.~\cite{chatterjee2017} explored GPU multiway mergesort
(MMS), employing multiway merge with partition-based decomposition.  While
effective at reducing merge passes, their approach requires complex
partitioning logic and does not exploit shared-memory merge trees.

\paragraph{ModernGPU.}
Baxter's ModernGPU library~\cite{moderngpu2013} provides high-performance
GPU merge routines built on the merge-path framework.  It serves as our
baseline for 2-way merge performance.

\paragraph{GPU External Sort.}
Stehle and Jacobsen~\cite{stehle2017} presented a GPU external sort at
SIGMOD 2017 that uses the GPU as a co-processor for merge operations,
streaming data between host and device.  Our external-sort pipeline follows
a similar streaming model but with higher merge fan-in.


% ══════════════════════════════════════════════════════════════════════════════
% 3  ARCHITECTURE
% ══════════════════════════════════════════════════════════════════════════════
\section{System Architecture}
\label{sec:architecture}

\gpucrocsort comprises three main components, illustrated in
\Cref{fig:architecture}: run generation, multi-pass merge, and external-sort
streaming.  We describe each in turn.

\begin{figure}[t]
  \placeholderfig{1.5cm}{System architecture diagram showing the three phases:
  (1) run generation via bitonic sort, (2) multi-pass merge with Strategy~A
  (2-way) or Strategy~B (8-way), and (3) external-sort streaming pipeline
  with pinned host memory.}
  \caption{Overall architecture of \gpucrocsort.}
  \label{fig:architecture}
\end{figure}

\subsection{Phase 1: Run Generation}
\label{sec:run_gen}

Run generation sorts small blocks of records entirely within GPU shared
memory using a \emph{bitonic sort} network.  Each thread block processes
512 records with 256 threads.  The procedure is as follows:

\begin{enumerate}[leftmargin=*,topsep=2pt,itemsep=1pt]
  \item Load 512 records from HBM into shared memory.
  \item Execute $\log_2 512 = 9$ stages of the bitonic sort network.
    Each stage consists of compare-and-swap operations that are
    naturally parallel across all 256 threads.
  \item Compute OVC delta values for adjacent records in the sorted
    output.  Because records are sorted, the OVC of each record
    relative to its predecessor can be computed by scanning the key
    from the offset where the keys first differ.
  \item Write the sorted run (records + OVC values) to HBM.
\end{enumerate}

For a dataset of $R$ records with block size $B{=}512$, this phase produces
$N = \lceil R/B \rceil$ sorted runs.

\begin{algorithm}[t]
  \caption{Bitonic Sort Run Generation}
  \label{alg:run_gen}
  \KwInput{Unsorted records $r[0..B{-}1]$ in HBM}
  \KwOutput{Sorted run with OVC values in HBM}
  \BlankLine
  \Kernel{BitonicSortKernel}{
    $\text{tid} \gets \text{threadIdx.x}$\;
    \Comment{Load into shared memory}
    $\text{smem}[\text{tid}] \gets r[\text{blockIdx.x} \cdot B + \text{tid}]$\;
    $\text{smem}[\text{tid} + 256] \gets r[\text{blockIdx.x} \cdot B + \text{tid} + 256]$\;
    $\text{\_\_syncthreads()}$\;
    \BlankLine
    \Comment{Bitonic sort network: 9 stages}
    \For{$\text{stage} \gets 0$ \KwTo $8$}{
      \For{$\text{step} \gets \text{stage}$ \KwTo $0$}{
        $\text{partner} \gets \text{tid} \oplus (1 \ll \text{step})$\;
        \If{$\text{partner} > \text{tid}$}{
          \text{CompareAndSwap}($\text{smem}[\text{tid}]$, $\text{smem}[\text{partner}]$)\;
        }
        $\text{\_\_syncthreads()}$\;
      }
    }
    \BlankLine
    \Comment{Compute OVC deltas}
    \If{$\text{tid} > 0$}{
      $\text{ovc}[\text{tid}] \gets \text{ComputeOVC}(\text{smem}[\text{tid}{-}1],\, \text{smem}[\text{tid}])$\;
    }
    \Else{
      $\text{ovc}[0] \gets \text{INITIAL\_OVC}$\;
    }
    $\text{\_\_syncthreads()}$\;
    \BlankLine
    \Comment{Write sorted run to HBM}
    $r[\text{blockIdx.x} \cdot B + \text{tid}] \gets \text{smem}[\text{tid}]$\;
    $\text{ovcOut}[\text{blockIdx.x} \cdot B + \text{tid}] \gets \text{ovc}[\text{tid}]$\;
  }
\end{algorithm}

\subsection{Phase 2: Multi-Pass Merge}
\label{sec:merge_phase}

Given $N$ sorted runs from Phase~1, the merge phase reduces them to a single
sorted output through iterative merge passes.  We implement two strategies:

\paragraph{Strategy A: 2-Way Merge Path (Baseline).}
Each pass merges pairs of adjacent runs using the merge-path
decomposition~\cite{green2012}.  All GPU threads participate in each merge
operation via balanced partitioning.  This strategy requires
$\lceil\log_2 N\rceil$ passes.

\paragraph{Strategy B: 8-Way Shared-Memory Merge Tree (Novel).}
Each pass merges groups of 8 runs using our shared-memory K-way merge tree
(\Cref{sec:kway}).  This strategy requires only $\lceil\log_8 N\rceil$
passes---a $3\times$ reduction---while maintaining full thread utilization
at every stage.

\subsection{Phase 3: External Sort Streaming}
\label{sec:external}

When the dataset exceeds GPU HBM capacity, \gpucrocsort employs a streaming
pipeline with pinned host memory:

\begin{enumerate}[leftmargin=*,topsep=2pt,itemsep=1pt]
  \item Partition the dataset into HBM-sized chunks on the host.
  \item Asynchronously transfer each chunk to the GPU using CUDA streams
    and pinned (\texttt{cudaHostAlloc}) memory.
  \item Execute run generation and as many merge passes as possible
    on-device.
  \item Stream partially merged results back to the host.
  \item Perform a final merge pass on partially merged runs, potentially
    multi-pass if necessary, using the GPU for each merge batch.
\end{enumerate}

Double-buffering with two CUDA streams hides transfer latency behind
compute, achieving near-peak HBM utilization.


% ══════════════════════════════════════════════════════════════════════════════
% 4  SHARED-MEMORY K-WAY MERGE TREE
% ══════════════════════════════════════════════════════════════════════════════
\section{Shared-Memory K-Way Merge Tree}
\label{sec:kway}

This section presents our primary contribution: a decomposition of $K$-way
merge into a tree of fully parallel 2-way merges executed entirely in GPU
shared memory.

\subsection{Motivation}

The fundamental tension in GPU merge sort is between \emph{fan-in} and
\emph{parallelism}.  \Cref{tab:tradeoff} summarizes the trade-off.

\begin{table}[t]
  \centering
  \caption{Fan-in vs.\ parallelism trade-off for 19{,}531 runs.}
  \label{tab:tradeoff}
  \begin{tabular}{@{}lccc@{}}
    \toprule
    \textbf{Strategy} & \textbf{Passes} & \textbf{Threads Active} & \textbf{HBM Traffic} \\
    \midrule
    2-way merge path  & 15 & All (256/block)     & $30\times$ data \\
    8-way loser tree   & 5  & 1 per merge         & $10\times$ data \\
    \textbf{8-way merge tree (ours)} & \textbf{5} & \textbf{All (256/block)} & $\mathbf{10\times}$ \textbf{data} \\
    \bottomrule
  \end{tabular}
\end{table}

The 2-way merge-path approach delivers full thread utilization but
wastes bandwidth with excessive passes.  The traditional K-way loser tree
minimizes passes but serializes the merge.  Our approach achieves
\emph{both} simultaneously.

\subsection{Tree Decomposition}
\label{sec:tree_decomp}

For $K{=}8$, the 8-way merge is decomposed into $\log_2 8 = 3$ levels of
2-way merges, organized as a binary tree (\Cref{fig:merge_tree}):

\begin{figure}[t]
  \placeholderfig{3cm}{Merge tree diagram showing 3 levels for $K{=}8$.
  Level~0: four 2-way merges (runs 0+1, 2+3, 4+5, 6+7), each with 64
  threads.  Level~1: two 2-way merges combining Level~0 outputs, each
  with 128 threads.  Level~2: one final 2-way merge with 256 threads.
  Arrows show data flow through shared-memory ping-pong buffers.  Only
  the Level~2 output writes to HBM.}
  \caption{Shared-memory 8-way merge tree.  Thread allocation per
  level is shown in parentheses.  All 256 threads are active at every
  level.}
  \label{fig:merge_tree}
\end{figure}

\paragraph{Level 0: Leaf Merges.}
Four independent 2-way merges, each handled by 64 threads using
merge-path decomposition.  Input: 8 sorted runs from HBM (or from the
previous pass's output).  Output: 4 merged sequences in shared-memory
buffer~A.

\paragraph{Level 1: Intermediate Merges.}
Two independent 2-way merges, each handled by 128 threads.  Input:
4 sequences from shared-memory buffer~A.  Output: 2 merged sequences in
shared-memory buffer~B (ping-pong).

\paragraph{Level 2: Root Merge.}
One 2-way merge handled by all 256 threads.  Input: 2 sequences from
shared-memory buffer~B.  Output: final merged sequence written to HBM.

\subsection{Shared-Memory Ping-Pong Buffering}
\label{sec:pingpong}

The tree uses \emph{ping-pong buffering} between two shared-memory regions.
At each level, the input is read from one buffer and the merged output is
written to the other.  This avoids read-write conflicts without requiring
additional synchronization beyond the standard \texttt{\_\_syncthreads()}
barrier between levels.

The total shared memory required per thread block is:

\begin{equation}
  \text{SMEM} = 2 \times K \times B_{\text{tile}} \times \text{sizeof}(\text{record})
  \label{eq:smem}
\end{equation}

\noindent where $B_{\text{tile}}$ is the tile size per run.  For $K{=}8$,
$B_{\text{tile}}{=}64$ records, and 100-byte records, this is
$2 \times 8 \times 64 \times 100 = 102{,}400$ bytes, well within the
48\,KB--164\,KB shared-memory budget of modern GPU architectures.

\subsection{Thread Allocation}

A key design choice is how to distribute threads across levels.
The principle is simple: \emph{all 256 threads participate at every level}.
At Level~$\ell$, there are $K/2^{\ell+1}$ independent merges, each assigned
$256 / (K/2^{\ell+1}) = 256 \cdot 2^{\ell+1} / K$ threads.  For $K{=}8$:

\begin{table}[t]
  \centering
  \caption{Thread allocation across merge tree levels ($K{=}8$, 256 threads).}
  \label{tab:thread_alloc}
  \begin{tabular}{@{}ccccc@{}}
    \toprule
    \textbf{Level} & \textbf{Merges} & \textbf{Threads/Merge} & \textbf{Total Threads} & \textbf{Buffer} \\
    \midrule
    0 & 4 & 64  & 256 & A $\to$ B \\
    1 & 2 & 128 & 256 & B $\to$ A \\
    2 & 1 & 256 & 256 & A $\to$ HBM \\
    \bottomrule
  \end{tabular}
\end{table}

All 256 threads are active at every level, achieving 100\% occupancy
within the thread block.

\subsection{Merge-Path at Each Level}

Within each 2-way merge at each level, we apply the merge-path
decomposition.  For a merge with $T$ threads merging two sequences of
length $L$, each thread computes its merge-path diagonal
$d = \text{rank} \cdot (2L/T)$ and performs a binary search to find the
partition point.  The thread then sequentially merges its assigned portion
($2L/T$ elements) into the output.

\begin{algorithm}[t]
  \caption{Shared-Memory 8-Way Merge Tree}
  \label{alg:kway_merge}
  \KwInput{$K{=}8$ sorted runs in HBM, each of length $L$}
  \KwOutput{Single merged run of length $8L$ in HBM}
  \BlankLine
  \Kernel{KWayMergeTree}{
    $\text{tid} \gets \text{threadIdx.x}$ \Comment{0..255}
    \BlankLine
    \Comment{Load 8 runs into shared memory buffer A}
    \text{CooperativeLoad}($\text{smem\_A}$, HBM, $8L$)\;
    $\text{\_\_syncthreads()}$\;
    \BlankLine
    \Comment{--- Level 0: 4 independent 2-way merges ---}
    $\text{merge\_id} \gets \text{tid} / 64$ \Comment{0..3}
    $\text{local\_tid} \gets \text{tid} \bmod 64$\;
    $\text{in} \gets \text{smem\_A} + \text{merge\_id} \cdot 2L$\;
    $\text{out} \gets \text{smem\_B} + \text{merge\_id} \cdot 2L$\;
    \text{MergePath2Way}($\text{in}$, $L$, $\text{in}{+}L$, $L$,
                         $\text{out}$, $\text{local\_tid}$, 64)\;
    $\text{\_\_syncthreads()}$\;
    \BlankLine
    \Comment{--- Level 1: 2 independent 2-way merges ---}
    $\text{merge\_id} \gets \text{tid} / 128$ \Comment{0..1}
    $\text{local\_tid} \gets \text{tid} \bmod 128$\;
    $\text{in} \gets \text{smem\_B} + \text{merge\_id} \cdot 4L$\;
    $\text{out} \gets \text{smem\_A} + \text{merge\_id} \cdot 4L$\;
    \text{MergePath2Way}($\text{in}$, $2L$, $\text{in}{+}2L$, $2L$,
                         $\text{out}$, $\text{local\_tid}$, 128)\;
    $\text{\_\_syncthreads()}$\;
    \BlankLine
    \Comment{--- Level 2: 1 final 2-way merge ---}
    $\text{in} \gets \text{smem\_A}$\;
    $\text{out} \gets \text{HBM\_output}$\;
    \text{MergePath2Way}($\text{in}$, $4L$, $\text{in}{+}4L$, $4L$,
                         $\text{out}$, $\text{tid}$, 256)\;
  }
\end{algorithm}

\subsection{Correctness}

The correctness of the merge tree follows from two properties:

\begin{enumerate}[leftmargin=*,topsep=2pt,itemsep=1pt]
  \item \textbf{Composability of 2-way merge.}  If $A$, $B$, $C$, $D$
    are sorted sequences, then $\text{merge}(\text{merge}(A, B),\,
    \text{merge}(C, D)) = \text{merge}(A, B, C, D)$.  The tree
    decomposition applies this recursively.
  \item \textbf{Correctness of merge-path.}  Each 2-way merge at each
    level is performed via the merge-path decomposition, whose correctness
    is established in~\cite{green2012}.
\end{enumerate}

The \texttt{\_\_syncthreads()} barrier between levels ensures that all
outputs from level $\ell$ are complete before level $\ell{+}1$ reads them.


% ══════════════════════════════════════════════════════════════════════════════
% 5  OVC ON GPU
% ══════════════════════════════════════════════════════════════════════════════
\section{Offset-Value Coding on GPU}
\label{sec:ovc}

\subsection{Packed Format}

The 32-bit OVC encoding (\Cref{eq:ovc_format}) maps naturally to GPU
registers.  Each OVC fits in a single \texttt{uint32\_t}, and comparison
reduces to a single unsigned integer comparison instruction
(\texttt{ISETP.LT.U32} in PTX).

\subsection{Flag Types}

Five flag values encode record status, ordered so that the packed
\texttt{uint32\_t} representation preserves sort order:

\begin{table}[t]
  \centering
  \caption{OVC flag types and their encoded values.}
  \label{tab:ovc_flags}
  \begin{tabular}{@{}clcl@{}}
    \toprule
    \textbf{Value} & \textbf{Flag} & \textbf{Bits} & \textbf{Semantics} \\
    \midrule
    0 & \texttt{EarlyFence} & \texttt{000} & End-of-run sentinel (sorts first) \\
    1 & \texttt{Duplicate}  & \texttt{001} & Identical to reference key \\
    2 & \texttt{Normal}     & \texttt{010} & Standard OVC comparison \\
    3 & \texttt{Initial}    & \texttt{011} & First record (full key compare) \\
    4 & \texttt{LateFence}  & \texttt{100} & End-of-input sentinel (sorts last) \\
    \bottomrule
  \end{tabular}
\end{table}

The encoding ensures that sentinel values sort to the correct positions
(fences at extremes) and that duplicate keys are identified without
accessing the full key.

\subsection{OVC-Aware Loser Tree}

For the traditional (single-threaded) K-way merge baseline, we implement
an OVC-aware loser tree in GPU registers.  The key optimization is
the \emph{duplicate shortcut}: when two OVC values have the
\texttt{Duplicate} flag, they are known to be equal without key access.
This reduces memory traffic during the merge, as duplicate keys need not
be fetched from HBM.

\begin{algorithm}[t]
  \caption{OVC Comparison}
  \label{alg:ovc_compare}
  \KwInput{Two OVC values $a$, $b$ as \texttt{uint32\_t}}
  \KwOutput{Ordering: $-1$, $0$, or $+1$}
  \BlankLine
  \Fn{OvcCompare($a$, $b$)}{
    \Comment{Natural uint32 ordering matches sort order}
    \uIf{$a < b$}{
      \Return $-1$\;
    }
    \uElseIf{$a > b$}{
      \Return $+1$\;
    }
    \Else{
      \Comment{Equal OVC: must compare full keys}
      \Return \text{CompareFullKeys}($\text{key}(a)$, $\text{key}(b)$)\;
    }
  }
\end{algorithm}

\subsection{OVC Delta Computation}

After each merge step, OVC deltas must be recomputed for the merged
output relative to the new predecessor.  This is performed cooperatively
by all threads in the block:

\begin{equation}
  \text{ovc}[i] = \text{ComputeOVC}(\text{out}[i{-}1],\, \text{out}[i])
  \quad \forall\, i \in [1, n)
  \label{eq:ovc_delta}
\end{equation}

The \texttt{ComputeOVC} function scans the key bytes starting from
offset~0, finds the first byte position $p$ where the keys differ, and
packs $(p, \text{key}[i][p])$ into the OVC format with a \texttt{Normal}
flag (or \texttt{Duplicate} if the keys are identical).


% ══════════════════════════════════════════════════════════════════════════════
% 6  THEORETICAL ANALYSIS
% ══════════════════════════════════════════════════════════════════════════════
\section{Theoretical Analysis}
\label{sec:analysis}

\subsection{HBM Traffic Model}

Each merge pass reads and writes the entire dataset once, producing
$2 \times D$ bytes of HBM traffic (where $D$ is the data size).  The total
HBM traffic is:

\begin{equation}
  T_{\text{HBM}} = 2 \cdot D \cdot \lceil\log_K N\rceil
  \label{eq:hbm_traffic}
\end{equation}

\noindent where $N$ is the number of initial sorted runs and $K$ is the
merge fan-in.

\subsection{Comparative Analysis}

\Cref{tab:traffic} compares the two strategies for a 1\,GB dataset with
$B{=}512$ records/block (producing $N{=}19{,}531$ runs) and a 10\,GB
dataset ($N{=}195{,}313$).

\begin{table}[t]
  \centering
  \caption{HBM traffic comparison for different merge strategies.}
  \label{tab:traffic}
  \begin{tabular}{@{}lcccc@{}}
    \toprule
    & \multicolumn{2}{c}{\textbf{1\,GB (19.5K runs)}} & \multicolumn{2}{c}{\textbf{10\,GB (195K runs)}} \\
    \cmidrule(lr){2-3} \cmidrule(lr){4-5}
    \textbf{Strategy} & \textbf{Passes} & \textbf{Traffic} & \textbf{Passes} & \textbf{Traffic} \\
    \midrule
    2-way ($K{=}2$)   & 15 & 30\,GB & 18 & 360\,GB \\
    4-way ($K{=}4$)   &  8 & 16\,GB & 9  & 180\,GB \\
    8-way ($K{=}8$)   &  5 & 10\,GB & 6  & 120\,GB \\
    16-way ($K{=}16$)  &  4 &  8\,GB & 5  & 100\,GB \\
    \midrule
    \textbf{Reduction} ($K{=}8$ vs $K{=}2$) & $3\times$ & $3\times$ & $3\times$ & $3\times$ \\
    \bottomrule
  \end{tabular}
\end{table}

\subsection{Shared Memory Overhead}

The shared-memory merge tree introduces internal data movement within
shared memory.  At each level $\ell$, the data is read from one buffer and
written to another, contributing shared-memory traffic but \emph{no}
additional HBM traffic.  The total shared-memory traffic per merge-tree
invocation is:

\begin{equation}
  T_{\text{SMEM}} = 2 \cdot K \cdot L \cdot \text{sizeof}(\text{record}) \cdot \log_2 K
  \label{eq:smem_traffic}
\end{equation}

Since shared memory bandwidth exceeds HBM bandwidth by an order of
magnitude on modern GPUs (e.g., ${\sim}19$\,TB/s aggregate shared-memory
bandwidth on A100 vs.\ 2\,TB/s HBM), this overhead is negligible.

\subsection{Roofline Analysis}

The merge operation is memory-bandwidth-bound: each element is compared
(1 integer comparison) and written (1 store), giving an arithmetic
intensity of approximately $1/\text{sizeof}(\text{record})$ FLOP/byte.
For 100-byte GenSort records, this is $0.01$ FLOP/byte, far below the
roofline knee of any modern GPU.  Thus, minimizing HBM passes directly
translates to proportional speedup.


% ══════════════════════════════════════════════════════════════════════════════
% 7  EXPERIMENTAL SETUP AND RESULTS
% ══════════════════════════════════════════════════════════════════════════════
\section{Experiments}
\label{sec:experiments}

\subsection{Setup}

\paragraph{Data.}
We use GenSort~\cite{gensort} 100-byte records (10-byte key + 90-byte
value), the standard benchmark for sort competitions.  Dataset sizes range
from 100K to 100M records (10\,MB to 10\,GB).

\paragraph{Hardware.}
Experiments are conducted on two GPU architectures:
\begin{itemize}[leftmargin=*,topsep=2pt,itemsep=1pt]
  \item \textbf{NVIDIA A100 (80\,GB)}: sm\_80, 2,039\,GB/s HBM2e,
    40\,MB L2, 164\,KB shared memory per SM.
  \item \textbf{NVIDIA H100 (80\,GB)}: sm\_90, 3,350\,GB/s HBM3,
    50\,MB L2, 228\,KB shared memory per SM.
\end{itemize}

\paragraph{Software.}
CUDA~12.x, compiled with \texttt{nvcc -O3}.  Baseline comparisons use
CUB~\cite{cub} \texttt{DeviceRadixSort} and Thrust
\texttt{thrust::sort}.

\paragraph{Experiment Suite.}
We execute nine experiment programs covering different aspects:
\begin{enumerate}[leftmargin=*,topsep=2pt,itemsep=1pt]
  \item Bottleneck analysis (run generation vs.\ merge)
  \item AoS vs.\ SoA memory layout comparison
  \item CPU vs.\ GPU throughput scaling
  \item CUB radix sort crossover point
  \item Kernel profiling (occupancy, bandwidth utilization)
  \item Edge cases (pre-sorted, reverse-sorted, all-duplicates)
  \item External sort with data exceeding HBM
  \item OVC encoding overhead measurement
  \item Merge fan-in scaling ($K{=}2,4,8,16$)
\end{enumerate}

\subsection{Throughput Scaling}

\begin{figure}[t]
  \placeholderfig{3cm}{Line chart: x-axis = number of records (log scale,
  100K to 100M), y-axis = throughput in M records/sec.  Lines for:
  (1) GPU CrocSort 8-way, (2) GPU CrocSort 2-way, (3) CUB RadixSort,
  (4) CPU CrocSort.  Expected: GPU 8-way fastest for large sizes,
  CUB competitive for smaller sizes, CPU slowest.}
  \caption{Sorting throughput vs.\ dataset size.}
  \label{fig:throughput}
\end{figure}

\todo{Populate with experimental data.}

\subsection{Merge Strategy Comparison}

\begin{table}[t]
  \centering
  \caption{Merge strategy comparison on 1\,GB dataset (10M records).}
  \label{tab:merge_comparison}
  \begin{tabular}{@{}lcccc@{}}
    \toprule
    \textbf{Strategy} & \textbf{Passes} & \textbf{HBM Traffic} & \textbf{Time (ms)} & \textbf{BW Util.} \\
    \midrule
    2-way merge path    & 15 & 30\,GB & \todo{--} & \todo{--} \\
    8-way loser tree    & 5  & 10\,GB & \todo{--} & \todo{--} \\
    8-way merge tree    & 5  & 10\,GB & \todo{--} & \todo{--} \\
    \midrule
    CUB RadixSort       & -- & --     & \todo{--} & \todo{--} \\
    \bottomrule
  \end{tabular}
\end{table}

\todo{Fill in timing and bandwidth utilization results.}

\subsection{CPU vs.\ GPU Comparison}

\begin{table}[t]
  \centering
  \caption{CPU vs.\ GPU sort performance on 10M records.}
  \label{tab:cpu_gpu}
  \begin{tabular}{@{}lccc@{}}
    \toprule
    \textbf{System} & \textbf{Time (ms)} & \textbf{Throughput (M rec/s)} & \textbf{Speedup} \\
    \midrule
    CPU CrocSort (1 thread)  & \todo{--} & \todo{--} & $1\times$ \\
    CPU CrocSort (16 threads) & \todo{--} & \todo{--} & \todo{--} \\
    GPU 2-way merge path     & \todo{--} & \todo{--} & \todo{--} \\
    GPU 8-way merge tree     & \todo{--} & \todo{--} & \todo{--} \\
    \bottomrule
  \end{tabular}
\end{table}

\todo{Fill in CPU vs.\ GPU comparison data.}

\subsection{Memory Layout Impact}

\begin{table}[t]
  \centering
  \caption{Impact of memory layout on merge throughput (10M records).}
  \label{tab:layout}
  \begin{tabular}{@{}lcc@{}}
    \toprule
    \textbf{Layout} & \textbf{Merge Time (ms)} & \textbf{Effective BW (GB/s)} \\
    \midrule
    Array of Structures (AoS) & \todo{--} & \todo{--} \\
    Structure of Arrays (SoA) & \todo{--} & \todo{--} \\
    Key-only (10B)            & \todo{--} & \todo{--} \\
    \bottomrule
  \end{tabular}
\end{table}

\todo{Fill in memory layout comparison data.}

\subsection{External Sort Performance}

\begin{figure}[t]
  \placeholderfig{3cm}{Bar chart: x-axis = dataset size (1GB, 5GB, 10GB,
  50GB, 100GB), y-axis = sort time in seconds.  Stacked bars showing
  time breakdown: host-to-device transfer, run generation, merge passes,
  device-to-host transfer.  Expected: merge dominates for large sizes,
  transfer overhead significant for external sort.}
  \caption{External sort time breakdown by phase.}
  \label{fig:external}
\end{figure}

\todo{Populate with external sort experiments.}


% ══════════════════════════════════════════════════════════════════════════════
% 8  RELATED WORK
% ══════════════════════════════════════════════════════════════════════════════
\section{Related Work}
\label{sec:related}

\paragraph{CPU External Sort.}
\crocsort~\cite{crocsort2026} established a new state of the art for CPU
external merge sort through its use of Offset-Value Coding and careful
cache-aware merge strategies.  Our work extends its core ideas---OVC
encoding and analytically driven merge fan-in---to the GPU domain, where
the memory hierarchy (HBM $\to$ shared memory $\to$ registers) differs
fundamentally from the CPU hierarchy (DRAM $\to$ L3 $\to$ L2 $\to$ L1
$\to$ registers).

\paragraph{Merge Path.}
Green et al.~\cite{green2012} introduced the merge-path framework for
parallel 2-way merge on GPUs.  Their key insight is that a 2-way merge can
be decomposed into independent sub-problems by tracing a path through the
merge matrix.  We build directly on this primitive, using it as the
building block at each level of our merge tree.

\paragraph{GPU Multiway Mergesort.}
Chatterjee et al.~\cite{chatterjee2017} proposed GPU Multiway Mergesort
(MMS), which partitions multiway merge among GPU threads.  Their approach
uses a different partitioning strategy: they find split points in all $K$
input sequences simultaneously, then assign tiles to thread blocks.
In contrast, our approach decomposes the K-way merge hierarchically,
avoiding the complexity of K-way partitioning while achieving full
parallelism through the merge tree.

\paragraph{GPU External Sort.}
Stehle and Jacobsen~\cite{stehle2017} presented a GPU-accelerated external
sort at SIGMOD 2017.  Their system uses the GPU for merge operations while
streaming data between host and device.  Our external sort pipeline follows
a similar model but achieves higher merge fan-in through the shared-memory
merge tree, reducing the number of host-device round trips.

\paragraph{ModernGPU.}
Baxter's ModernGPU library~\cite{moderngpu2013} provides optimized
implementations of merge-path-based sorting and merging on CUDA GPUs.
It serves as both an inspiration and a baseline for our 2-way merge
implementation.

\paragraph{CUB and Thrust.}
The CUB~\cite{cub} and Thrust libraries provide GPU radix sort, which is
asymptotically optimal for fixed-width keys.  For GenSort's 10-byte keys,
radix sort requires 10 passes over the key bytes.  Our comparison-based
approach may be competitive for longer keys or variable-length keys where
radix sort's pass count grows proportionally.


% ══════════════════════════════════════════════════════════════════════════════
% 9  FUTURE WORK
% ══════════════════════════════════════════════════════════════════════════════
\section{Future Work}
\label{sec:future}

Several directions offer promising extensions of \gpucrocsort:

\paragraph{Higher Fan-In ($K{=}16, 32$).}
Increasing the merge fan-in to $K{=}16$ or $K{=}32$ would further reduce
pass counts ($\lceil\log_{16} N\rceil$ or $\lceil\log_{32} N\rceil$).
The challenge is shared-memory capacity: $K{=}16$ with 100-byte records
requires $2 \times 16 \times 64 \times 100 = 204{,}800$ bytes, which
exceeds the 164\,KB shared-memory limit on A100 but fits within the
228\,KB budget on H100.

\paragraph{Key-Value Separation.}
Separating keys from values in the merge pipeline would reduce data
movement during the merge tree levels, where only keys (and OVC values)
need to participate in comparisons.  Values would be gathered only at
the final output stage, reducing shared-memory pressure by up to
$9\times$ for GenSort records.

\paragraph{Multi-GPU Scaling.}
NCCL-based multi-GPU external sort, where each GPU sorts a partition and
a final merge phase combines results across GPUs via NVLink or
NVSwitch interconnect.

\paragraph{GPU Direct Storage.}
Bypassing the CPU entirely for external sort by using GPU Direct Storage
(GDS) to read and write sorted runs directly between NVMe SSDs and GPU
HBM, eliminating PCIe round trips through host memory.

\paragraph{OVC-Aware Merge Path.}
Incorporating OVC values into the merge-path binary search to exploit
the comparison shortcut.  When OVC values indicate that one sequence's
elements are all less than the other's (detectable via fence values),
the binary search can short-circuit.

\paragraph{H100-Specific Optimizations.}
Exploiting H100's architectural features:
\begin{itemize}[leftmargin=*,topsep=2pt,itemsep=1pt]
  \item \textbf{SM Clusters:} Distributed shared memory across clustered
    SMs could enable larger merge tiles without spilling to HBM.
  \item \textbf{Tensor Memory Accelerator (TMA):} Hardware-accelerated
    asynchronous data movement between HBM and shared memory.
  \item \textbf{\texttt{cp.async}:} Asynchronous copy instructions to
    overlap data loading with merge computation.
\end{itemize}


% ══════════════════════════════════════════════════════════════════════════════
% 10  CONCLUSION
% ══════════════════════════════════════════════════════════════════════════════
\section{Conclusion}
\label{sec:conclusion}

We have presented \gpucrocsort, a GPU-accelerated external merge sort
that resolves the fundamental tension between merge fan-in and thread
utilization through a novel shared-memory K-way merge tree.  By
decomposing a $K$-way merge ($K{=}8$) into $\log_2 K = 3$ levels of
fully parallel 2-way merge-path operations within shared memory, our
approach achieves the pass-count reduction of K-way merging
($\lceil\log_8 N\rceil$ passes instead of $\lceil\log_2 N\rceil$)
while maintaining 100\% thread utilization at every level.

For a 1\,GB dataset producing 19{,}531 sorted runs, this translates to a
$3\times$ reduction in HBM traffic (10\,GB vs.\ 30\,GB) compared to the
2-way merge-path baseline.  Combined with GPU-adapted Offset-Value
Coding for single-instruction key comparison and a streaming external-sort
pipeline for datasets exceeding HBM capacity, \gpucrocsort demonstrates
that the GPU memory hierarchy---with its high-bandwidth HBM and
ultra-fast shared memory---is well suited to external merge sort when
the merge algorithm is carefully matched to the hardware.

The shared-memory merge tree is a general technique applicable beyond
sorting: any hierarchical reduction where intermediate results fit in
shared memory can benefit from this decomposition.  We believe this
approach opens new avenues for GPU-accelerated data processing algorithms
that must balance fan-in with parallelism.


% ══════════════════════════════════════════════════════════════════════════════
% ACKNOWLEDGMENTS
% ══════════════════════════════════════════════════════════════════════════════
\section*{Acknowledgments}

\todo{Add acknowledgments.}


% ══════════════════════════════════════════════════════════════════════════════
% REFERENCES
% ══════════════════════════════════════════════════════════════════════════════
\bibliographystyle{IEEEtran}

\begin{thebibliography}{10}

\bibitem{crocsort2026}
CrocSort: External Merge Sort with Offset-Value Coding.
\newblock In \emph{Proceedings of the VLDB Endowment}, 2026.

\bibitem{ovc_original}
P.~Larson and G.~Graefe.
\newblock Offset-value coding.
\newblock In \emph{IEEE International Conference on Data Engineering (ICDE)}, 2014.

\bibitem{green2012}
O.~Green, R.~McColl, and D.~A.~Bader.
\newblock GPU merge path: A GPU merging algorithm.
\newblock In \emph{Proceedings of the 26th ACM International Conference on
  Supercomputing (ICS)}, pages 331--340, 2012.

\bibitem{chatterjee2017}
S.~Chatterjee, N.~Rai, and J.~D.~Owens.
\newblock GPU multiway mergesort.
\newblock In \emph{ACM SIGPLAN Symposium on Principles and Practice of
  Parallel Programming (PPoPP)}, 2017.

\bibitem{moderngpu2013}
S.~Baxter.
\newblock ModernGPU: Patterns and techniques for GPU computing.
\newblock \url{https://moderngpu.github.io/}, 2013.

\bibitem{stehle2017}
E.~Stehle and H.-A.~Jacobsen.
\newblock A memory bandwidth-efficient hybrid radix sort on GPUs.
\newblock In \emph{Proceedings of the ACM SIGMOD International Conference on
  Management of Data}, pages 417--432, 2017.

\bibitem{cub}
D.~Merrill.
\newblock CUB: CUDA unbound.
\newblock \url{https://nvlabs.github.io/cub/}, 2015.

\bibitem{gensort}
C.~Nyberg.
\newblock GenSort data generator.
\newblock \url{http://www.ordinal.com/gensort.html}.

\bibitem{cuda_programming_guide}
NVIDIA.
\newblock \emph{CUDA C++ Programming Guide}, v12.x.
\newblock NVIDIA Corporation, 2024.

\bibitem{nsight}
NVIDIA.
\newblock \emph{Nsight Compute Profiling Guide}.
\newblock NVIDIA Corporation, 2024.

\end{thebibliography}

\end{document}
